\begin{abstract}
Robot learning now advances along two tracks that rarely meet. On one side, direct
Vision-Language-Action (VLA) policies map observations to actions and are scored by closed-loop task
success. On the other, predictive and generative world models forecast future observations and are
scored by open-loop prediction or generation quality. A natural question sits between them: does
world modelling earn a measurable, closed-loop \emph{advantage} over a direct policy, and for which
robotic capabilities? We argue that the field cannot yet answer this question, and that the reason
is a gap in how it is measured, not in the models themselves. World-model benchmarks score
prediction without ever executing it, while task-success suites host a single policy and never build
a world-model versus VLA contrast.

This survey maps the evaluation landscape around that gap. We catalogue 160 web-verified benchmarks
spanning 2017 to 2026 and organise them by evaluation mode, robotic capability, and model family
into four lanes: policy suites, embodied agents, world-model evaluation, and prediction-to-action
bridges. Across the corpus, 138 of 160 benchmarks are model-agnostic and only 11 (7\%) build an
explicit VLA-versus-world-model contrast; counterfactual capability is almost entirely unmeasured,
and only four benchmarks turn prediction into executed action. We contribute an operational
taxonomy, a coverage comparison against the eight closest surveys (ours is the only one to cross
capability with model family), an evaluation loop that isolates the advantage of prediction, and an
actionable protocol of four advantage-aware metrics anchored on named testbeds. The organising claim
is not that world models help or do not help, but that answering the question requires benchmarks
built to ask it.
\end{abstract}
